% Created 2025-11-14 Fri 14:03
% Intended LaTeX compiler: pdflatex
\documentclass[a4paper]{article}
\usepackage[utf8]{inputenc}
\usepackage[T1]{fontenc}
\usepackage{graphicx}
\usepackage{longtable}
\usepackage{wrapfig}
\usepackage{rotating}
\usepackage[normalem]{ulem}
\usepackage{amsmath}
\usepackage{amssymb}
\usepackage{capt-of}
\usepackage{hyperref}
\usepackage[spanish]{babel}
\usepackage[usenames,dvipsnames]{color} % Required for custom colors
\renewcommand{\ttdefault}{pcr} % MONOESPACIO CON NEGRIT
\usepackage{lastpage}
\usepackage{listings}
\usepackage{listingsutf8}
\renewcommand{\lstlistingname}{Listado}
\lstset{frame=single,inputencoding=utf8,basicstyle=\scriptsize\ttfamily,showstringspaces=false,numbers=none}
\definecolor{MyDarkGreen}{rgb}{0.0,0.4,0.0} % This is the color used for comments
\lstset{ breaklines=true, postbreak=\mbox{\textcolor{red}{$\hookrightarrow$}\space}, keywordstyle=\bfseries, keywordstyle=[1]\color{Blue}\bfseries,  keywordstyle=[2]\color{Purple}\bfseries,  keywordstyle=[3]\color{Blue}\underbar,   identifierstyle=,   commentstyle=\usefont{T1}{pcr}{m}{sl}\color{MyDarkGreen}\small,   stringstyle=\color{Purple},   showstringspaces=false,   tabsize=2,   morecomment=[l][\color{Blue}]{...} }
\lstset{literate=  {á}{{\'a}}1 {é}{{\'e}}1 {í}{{\'i}}1 {ó}{{\'o}}1 {ú}{{\'u}}1   {Á}{{\'A}}1 {É}{{\'E}}1 {Í}{{\'I}}1 {Ó}{{\'O}}1 {Ú}{{\'U}}1   {à}{{\`a}}1 {è}{{\`e}}1 {ì}{{\`i}}1 {ò}{{\`o}}1 {ù}{{\`u}}1   {À}{{\`A}}1 {È}{{\'E}}1 {Ì}{{\`I}}1 {Ò}{{\`O}}1 {Ù}{{\`U}}1   {ä}{{\"a}}1 {ë}{{\"e}}1 {ï}{{\"i}}1 {ö}{{\"o}}1 {ü}{{\"u}}1   {Ä}{{\"A}}1 {Ë}{{\"E}}1 {Ï}{{\"I}}1 {Ö}{{\"O}}1 {Ü}{{\"U}}1   {â}{{\^a}}1 {ê}{{\^e}}1 {î}{{\^i}}1 {ô}{{\^o}}1 {û}{{\^u}}1   {Â}{{\^A}}1 {Ê}{{\^E}}1 {Î}{{\^I}}1 {Ô}{{\^O}}1 {Û}{{\^U}}1   {œ}{{\oe}}1 {Œ}{{\OE}}1 {æ}{{\ae}}1 {Æ}{{\AE}}1 {ß}{{\ss}}1   {ű}{{\H{u}}}1 {Ű}{{\H{U}}}1 {ő}{{\H{o}}}1 {Ő}{{\H{O}}}1   {ç}{{\c c}}1 {Ç}{{\c C}}1 {ø}{{\o}}1 {å}{{\r a}}1 {Å}{{\r A}}1   {€}{{\euro}}1 {£}{{\pounds}}1 {«}{{\guillemotleft}}1   {»}{{\guillemotright}}1 {ñ}{{\~n}}1 {Ñ}{{\~N}}1 {¿}{{?`}}1 }
\usepackage{caption}
\usepackage{attachfile2}
\usepackage[margin=1.5cm,includeheadfoot,includehead,includefoot]{geometry}
\hypersetup{colorlinks,linkcolor=black}
\usepackage{fancyhdr}
\pagestyle{fancyplain}
\chead{}
\lhead{}
\rhead{}
\cfoot{}
\lfoot{\begin{footnotesize}alvaro.gonzalezsotillo@educa.madrid.org\end{footnotesize}}
\rfoot{\begin{footnotesize}\thepage / \pageref{LastPage}\end{footnotesize}}
\usepackage{svg}
\usepackage{letltxmacro}
\LetLtxMacro{\originalincludegraphics}{\includegraphics}
\renewcommand{\includegraphics}[2][]{\IfFileExists{#2.pdf}{\originalincludegraphics[#1]{#2.pdf}}{\originalincludegraphics[#1]{#2}}}
\LetLtxMacro{\originalincludesvg}{\includesvg}
\renewcommand{\includesvg}[2][]{\IfFileExists{#2.pdf}{\originalincludegraphics[#1]{#2.pdf}}{\originalincludegraphics[#1]{#2.svg.pdf}}}
\usepackage{comment}
\excludecomment{NOTES}
\author{Álvaro González Sotillo}
\date{\today}
\title{Elementos del desarrollo de Software}
\hypersetup{
 pdfauthor={Álvaro González Sotillo},
 pdftitle={Elementos del desarrollo de Software},
 pdfkeywords={},
 pdfsubject={},
 pdfcreator={Emacs 30.2 (Org mode 9.7.11)}, 
 pdflang={Spanish}}
\begin{document}

\maketitle
\setcounter{tocdepth}{1}
\tableofcontents

\captionsetup{font=scriptsize}
\section{IDE}
\label{sec:org0000018}
\begin{itemize}
\item \textbf{I} ntegrated \textbf{D} evelopment \textbf{E} nvironment
\item Una aplicación con conjunto de herramientas que facilitan la tarea al programador
\begin{itemize}
\item Editor de texto
\item Compilador
\item Enlazador
\item Control de versiones
\item Depurador
\item Integración con IA generativa
\item Empaquetado, generador de instalaciones
\end{itemize}
\end{itemize}
\subsection{Editor de código}
\label{sec:org0000000}
Proporciona una interfaz para escribir y editar el código fuente. Suele incluir características como resaltado de sintaxis, autocompletado, indentación automática y otras utilidades para facilitar la escritura del código.
\subsection{Compilador/Intérprete}
\label{sec:org0000003}
Permite compilar o interpretar el código fuente en un formato ejecutable o en bytecode, dependiendo del lenguaje de programación utilizado. Algunos IDEs también ofrecen la capacidad de ejecutar el código directamente desde el entorno.
\subsection{Depurador}
\label{sec:org0000006}
Es una herramienta que ayuda a identificar y corregir errores en el código. Permite establecer puntos de interrupción, examinar el estado de las variables, ejecutar el código paso a paso y realizar otras operaciones para analizar y solucionar problemas en el programa.
\subsection{Gestión de proyectos/ficheros}
\label{sec:org0000009}
Permite crear, organizar y administrar proyectos de desarrollo de software. Esto incluye la capacidad de crear estructuras de directorios, agregar o eliminar archivos, gestionar dependencias y realizar otras tareas relacionadas con la organización del proyecto.
\subsection{Control de versiones}
\label{sec:org000000c}
Algunos IDEs incluyen integración con sistemas de control de versiones como Git, que permiten realizar seguimiento de cambios en el código, realizar confirmaciones (commits), fusionar (merge) ramas y otras operaciones relacionadas con la gestión del código fuente.
\subsection{Herramientas de construcción}
\label{sec:org000000f}
Algunos IDEs proporcionan herramientas para automatizar el proceso de construcción del software, como la generación de archivos de configuración, la compilación, el empaquetado y otras tareas relacionadas con la construcción del proyecto.
\subsection{Herramientas externas}
\label{sec:org0000012}
Los IDEs suelen ofrecer integración con otras herramientas y servicios externos, como sistemas de gestión de bases de datos, terminal, servidores web, frameworks, bibliotecas, entre otros, para facilitar el desarrollo y la integración con otros componentes del sistema.
\subsection{\emph{Agentic code}}
\label{sec:org0000015}
Integración con IA generativa, que permite realizar funciones que antes se realizaban manualmente o con plugins específicos
\begin{itemize}
\item Generación de código
\item Modificación de código (\emph{refactor})
\end{itemize}
\section{Plugins}
\label{sec:org000001b}
\begin{itemize}
\item Complementos que se relacionan con otras herramientas para agregarle una nueva función.
\item Esta aplicación adicional es ejecutada por la aplicación principal.
\end{itemize}
\section{Proyecto de ejemplo}
\label{sec:org0000024}
\begin{itemize}
\item \href{sin-construccion.zip}{Proyecto de ejemplo}
\item Descomprime el ZIP
\end{itemize}
\subsection{Construcción sin IDE}
\label{sec:org000001e}
\begin{itemize}
\item Carpeta \texttt{Tetris}
\item Fichero \texttt{Makefile}
\begin{itemize}
\item Es un sistema básico de construcción
\item Antiguo, pero muy estándar y muy utilizado
\end{itemize}
\end{itemize}

\begin{lstlisting}[language=make,numbers=none]
src=src/*.java
pkg=tetris.jar
main=PlayTetris


all:
	mkdir -p bin
	javac ${src} -d bin/

package: all
	../package.sh ${pkg} ${main}

doc:
	javadoc ${src} -d doc/

play: all
	java -cp bin/ ${main}

clean:
	rm -rf bin
	rm -rf doc
	rm -f ${pkg}
\end{lstlisting}
\subsection{Ejercicio}
\label{sec:org0000021}
\begin{itemize}
\item Utiliza solo la línea de comandos
\item Construye la aplicación con los comandos que se deducen del fichero \texttt{Makefile}
\item Ejecuta la aplicación
\item Extra: utiliza el comando \texttt{make}
\end{itemize}
\section{Intellij}
\label{sec:org000002d}
\begin{itemize}
\item IDE profesional
\begin{itemize}
\item Originalmente para Java
\item Pero soporta otros lenguajes
\end{itemize}
\end{itemize}
\subsection{IDEs basados en IntelliJ}
\label{sec:org0000027}

\begin{center}
\begin{tabular}{llll}
IDE & Enfoque Principal & Características Clave & Audiencia Objetivo\\
\hline
\textbf{IntelliJ IDEA Ultimate} & Desarrollo Java y web empresarial & - IDE más completo, Soporte total para desarrollo empresarial & Desarrolladores profesionales, equipos empresariales\\
\textbf{IntelliJ IDEA Community Edition} & Desarrollo Java de código abierto & - Versión gratuita, Características básicas de Java & Desarrolladores individuales, estudiantes\\
\textbf{PyCharm} & Desarrollo en Python & - Entorno especializado para Python, Herramientas de análisis de datos & Programadores Python, científicos de datos\\
\textbf{WebStorm} & Desarrollo web y JavaScript & - Soporte avanzado para JavaScript Herramientas de frontend & Desarrolladores web front-end y full-stack\\
\textbf{CLion} & Desarrollo en C y C++ & - Soporte multiplataforma- Herramientas de optimización & Programadores de sistemas, desarrollo de juegos\\
\textbf{GoLand} & Desarrollo en Go & - Soporte para microservicios & Desarrolladores de Go, ingenieros backend\\
\textbf{Rider} & Desarrollo .NET &  & Desarrolladores .NET y C\#\\
\textbf{Android Studio} & Desarrollo de aplicaciones Android & - Herramienta oficial de Google, Soporte nativo para Kotlin/Java, Emulador integrado & Desarrolladores de aplicaciones móviles\\
\end{tabular}
\end{center}
\subsection{Por qué tantas versiones}
\label{sec:org000002a}
\begin{itemize}
\item La aplicación IntelliJ es prácticamente la misma
\item Las diferentes versiones se diferencian en los \emph{plugins} instalados
\item No todos los plugins se pueden instalar en todas las versiones
\begin{itemize}
\item La razón no es técnica, sino comercial
\end{itemize}
\end{itemize}
\section{Instalación Intellij}
\label{sec:org0000030}
\begin{itemize}
\item Como (casi) todos los IDE
\begin{itemize}
\item \href{https://www.jetbrains.com/idea/download/?section=linux}{Bajar la instalación} (la \emph{Comunity Edition}), o en el NAS
\item Descomprimir la instalación
\item Ejecutar \texttt{bin/idea.sh/}
\item Opcional: crear fichero \texttt{.desktop} para lanzarlo cómodamente
\end{itemize}
\item Algunos IDE complejos necesitan una instalación real
\begin{itemize}
\item Pero se debe más al S.O. que al IDE
\end{itemize}
\end{itemize}
\section[Conceptos Intellij]{Conceptos Intellij\hfill{}\textsc{ATTACH}}
\label{sec:org000003f}
\begin{itemize}
\item Proyecto
\item Módulo
\item Librería
\begin{itemize}
\item Fichero \texttt{jar}: clases precompiladas de Java
\item SDK (compilador del lenguaje). En Java, se llama JDK
\end{itemize}
\item Facet: capacidades que se añaden a un proyecto
\begin{itemize}
\item Ejemplo: añadir un lenguaje
\end{itemize}
\item Artifacts: creación de ficheros entregables/desplegables/ejecutables      

\begin{center}
\includegraphics[width=.9\linewidth]{/home/alvaro/apuntes-clase/entornos-de-desarrollo/apuntes/2/data/d0/909688-03e4-4791-8c4b-2f2fcac11fda/clipboard-20251113T192607.png}
\end{center}
\end{itemize}
\subsection{Proyecto}
\label{sec:org0000033}
\begin{itemize}
\item Un proyecto contiene código fuente y otros recursos para construir una aplicación
\item La aplicación se compone de módulos
\begin{itemize}
\item Debe tener al menos un módulo
\end{itemize}
\end{itemize}
\subsection{Módulo}
\label{sec:org0000036}
\begin{itemize}
\item Varios ficheros fuente y recursos en una unidad compilable
\item Un módulo se compila completo o no puede ejecutarse
\end{itemize}
\subsection{Libería}
\label{sec:org0000039}
\begin{itemize}
\item Código java ya compilado
\begin{itemize}
\item Un fichero \texttt{jar} ya generado
\item Otro módulo del proyecto
\end{itemize}
\item Se pueden añadir a cada módulo
\end{itemize}
\subsection{Artifact/entregable}
\label{sec:org000003c}
\begin{itemize}
\item Un fichero que es el resultado de la aplicación
\begin{itemize}
\item Por ejemplo, un \texttt{jar} ejecutable
\end{itemize}
\end{itemize}
\section{Ejercicio}
\label{sec:org0000042}
\begin{itemize}
\item Crea un módulo para cada juego
\item Compila y juega al \texttt{tetris}
\begin{itemize}
\item y al \texttt{mathhero}
\end{itemize}
\item Opcional: juega con otro compañero al \texttt{pongserver}
\end{itemize}
\section{Ejercicio}
\label{sec:org0000045}
\begin{itemize}
\item Crea un módulo nuevo de nombre \texttt{all}
\item Mostrará un menú modo texto para elegir el juego a lanzar
\item Lanzará el juego y al acabar el juego se acabará el programa
\end{itemize}
\end{document}
